%============================================================
% HALAMAN JUDUL DALAM BUKU KP
% Header/footer ditangani oleh pagestyle kpruler.
% Mendukung 1, 2, atau 3 mahasiswa.
%============================================================
\thispagestyle{empty}

\begin{center}
  \vspace{0.3cm}

  {\FontJudulDalamLabel\bfseries LAPORAN KERJA PRAKTIK\par}
  \vspace{1cm}

  \begin{minipage}{0.92\textwidth}
    \centering
    {\FontJudulDalamJudul\bfseries\MakeUppercase{\JudulKP}\par}
  \end{minipage}
  
  \vspace{0.8cm}
  {\FontJudulDalamIsi\bfseries Periode \PeriodeKP\par}

  \vfill

  \includegraphics[height=5.5cm]{assets/pens/logo_pens.png}\par

  \vfill

  {\FontJudulDalamIsi\bfseries Oleh:\par}
  \vspace{0.3cm}

  \begin{tabular}{@{}p{5.5cm}l@{}}
    \bfseries\FontJudulDalamIsi \NamaMahasiswaSatu & \bfseries\FontJudulDalamIsi NRP. \NRPSatu \\
    \ifcsvoid{NamaMahasiswaDua}{}{%
      \bfseries\FontJudulDalamIsi \NamaMahasiswaDua & \bfseries\FontJudulDalamIsi NRP. \NRPDua \\
    }
    \ifcsvoid{NamaMahasiswaTiga}{}{%
      \bfseries\FontJudulDalamIsi \NamaMahasiswaTiga & \bfseries\FontJudulDalamIsi NRP. \NRPTiga \\
    }
  \end{tabular}

  \vfill

  {\FontJudulDalamIsi\bfseries \MakeUppercase{\JenjangProdi} \MakeUppercase{\NamaProdi}\par}
  {\FontJudulDalamIsi\bfseries \MakeUppercase{\NamaDepartemen}\par}
  {\FontJudulDalamIsi\bfseries POLITEKNIK ELEKTRONIKA NEGERI SURABAYA\par}
  {\FontJudulDalamIsi\bfseries \TahunKP\par}

\end{center}
\thispagestyle{empty} % tanpa nomor halaman (Pasal 4.2.2 no. 4b)
